\section{Related Work}
\label{sec:related}

\paragraph{Prompt compression.}
One family reduces prompt length by scoring and removing tokens or sentences. \textsc{LLMLingua}~\cite{llmlingua} and \textsc{LongLLMLingua}~\cite{longllmlingua}
use a small language model to estimate token informativeness; \textsc{LLMLingua-2}~\cite{llmlingua2}
distils the decision into a token classifier for task-agnostic compression;
Selective Context~\cite{selective_context} drops spans that carry little
self-information; Provence~\cite{provence} prunes sentences jointly with question
relevance; CPC~\cite{cpc} scores sentences with a context-aware encoder;
Quito-X~\cite{quitox} casts compression as an information-bottleneck problem and
approximates mutual information with cross-attention; DAC~\cite{dac} adds an
attention-aware term to an entropy criterion; the Perception
Compressor~\cite{perception_compressor} retrieves relevant demonstrations with guiding
questions before compressing at the token level. These methods differ mainly in how they
estimate relevance and redundancy. They share one assumption: that the prompt is ordinary
text whose only cost is length. None of them alters a value the model is allowed to read,
and none is evaluated on questions whose answer \emph{is} a protected value. The omission is
consequential. The strongest of these methods leaves $98.1\%$ of protected values
recoverable in our benchmark, and it preserves the answer for a query that names a sensitive
entity far less reliably than it preserves an ordinary answer span.

\paragraph{Text sanitization for inference.}
A second family rewrites protected spans before the prompt leaves the user side.
DYNTEXT~\cite{dyntext} substitutes sensitive tokens with semantically similar non-sensitive
neighbours drawn from a dynamic adjacency list, so that a fixed substitution table cannot be
attacked; PromptGraph~\cite{promptgraph} guides the rewrite with a dependency graph;
INTACT~\cite{intact} replaces identifiers with generalised candidates selected for
resistance to LLM inference attacks. Masking wrappers around published compressors belong
here too. These methods decide per span what to hide and treat utility as whatever survives
the rewrite. On our benchmark that framing has a measurable cost: a masking pipeline that
achieves zero leakage scores $0.1712$ on the $2{,}728$-case extension, because the span it
removes is the referent the question needs. Differential privacy~\cite{dp} and
$k$-anonymity~\cite{kanon} supply formal guarantees, but neither offers the per-query
referential behaviour this task requires; Section~\ref{sec:discussion} returns to that point.

\paragraph{Long-context question answering.}
LongBench~\cite{longbench} supplies the passages from which we derive the public-benchmark
variants, and SQuAD~v2~\cite{squad2} supplies the scoring function. Our contribution is
orthogonal to the retrieval and long-context modelling literature: we do not modify the
reader, only the text it receives, and we measure the effect at a fixed token budget.

\paragraph{Where this study sits.}
The two families above optimise different objectives, length and concealment, and each
reports utility on prompts whose answer is ordinary text. We study the case where both
objectives fall on the same span, and show that the resulting conflict is resolvable by
changing \emph{which operation} acts on that span rather than by trading one objective
against the other. We keep the compression interface of the first family and the
per-span rewriting of the second, and modify the stage at which the operation is chosen.
